% GENERATED FILE. Edit canonical lessons, then run npm run generate:books.
% canonical-source-hash: fnv1a64:ce4637e132c38b14
% canonical-lessons: MR-C360-samanta, MR-C360-bhedbhav, MR-C360-visamta, MR-C360-daridrya, MR-C360-saksarta

\chapter{Fair and Unfair}
\label{ch:mr-fair-and-unfair}
\hlchaptermodality{360}

\begin{chapteropening}
\textbf{By the end of this chapter:} \emph{I can talk about equality, discrimination and poverty.}

Complete the last lesson of chapter 360: I can talk about equality, discrimination and poverty.
\end{chapteropening}

\section[\texorpdfstring{\mr{समानता} (samāntā)}{समानता (samāntā)}]{\mr{समानता} (samāntā) — equality}

\label{lesson:MR-C360-samanta}



\begin{quote}
Before the new one: say the Marathi for a verdict, then the Marathi for a claim.
\end{quote}

\subsection*{You'll want to know: \mr{समानता}}
\textbf{\mr{समानता}} — \emph{samāntā} — \textquotedblleft{}equality\textquotedblright{}.

\begin{quote}
\mr{स्त्री} \mr{आणि} \mr{पुरुषांमध्ये} \mr{समानता} \mr{असली} \mr{पाहिजे}.
\end{quote}

\begin{quote}
\emph{strī āṇi puruṣāmmadhye samāntā aslī pāhije.}
\end{quote}

\begin{quote}
There should be equality between women and men.
\end{quote}

Feminine.

\mr{समान}, equal, + -\mr{ता}, -ness.

\begin{cousinweb}
\textbf{\mr{समानता}} (\emph{samāntā}) is a tatsama, taken from Sanskrit \emph{samānatā}, \textquotedblleft{}equality, sameness\textquotedblright{} (Monier-Williams).
\end{cousinweb}

\begin{grammarlens}[title={What you've built}]
One more word for giving reasons and opinions.
\end{grammarlens}

\subsection*{Your turn}
\begin{itemize}
\raggedright
  \item \emph{Say it:} \emph{samāntā}
  \item \emph{Say it:} \emph{samāntā}, once more
  \item \emph{Say it:} \emph{strī āṇi puruṣāmmadhye samāntā aslī pāhije.} — There should be equality between women and men.
  \item \emph{Recall:} say the Marathi for a verdict, then the Marathi for a claim, then say \emph{samāntā} again
\end{itemize}

\begin{tcolorbox}[breakable,colback=teal!4,colframe=teal!35!black,title={Before you move on}]
What does \mr{समानता} mean? (\textquotedblleft{}Equality\textquotedblright{}.) Say it once more.
\end{tcolorbox}
\section[\texorpdfstring{\mr{भेदभाव} (bhedbhāv)}{भेदभाव (bhedbhāv)}]{\mr{भेदभाव} (bhedbhāv) — discrimination}

\label{lesson:MR-C360-bhedbhav}



\begin{quote}
Before the new one: say the Marathi for a claim, then the Marathi for equality.
\end{quote}

\subsection*{You'll want to know: \mr{भेदभाव}}
\textbf{\mr{भेदभाव}} — \emph{bhedbhāv} — \textquotedblleft{}discrimination\textquotedblright{}.

\begin{quote}
\mr{आमचे} \mr{शिक्षक} \mr{कोणताही} \mr{भेदभाव} \mr{करत} \mr{नाहीत}.
\end{quote}

\begin{quote}
\emph{āmce śikṣak koṇtāhī bhedbhāv karat nāhīt.}
\end{quote}

\begin{quote}
Our teachers show no discrimination at all.
\end{quote}

Masculine.

\mr{भेदभाव} \mr{करणे}, to discriminate.

\begin{cousinweb}
\textbf{\mr{भेदभाव}} (\emph{bhedbhāv}) is a learned compound built from Sanskrit \emph{bheda}, \textquotedblleft{}splitting, difference\textquotedblright{}, and \emph{bhāva}, \textquotedblleft{}state, attitude\textquotedblright{} (Monier-Williams).
\end{cousinweb}

\begin{grammarlens}[title={What you've built}]
One more word for giving reasons and opinions.
\end{grammarlens}

\subsection*{Your turn}
\begin{itemize}
\raggedright
  \item \emph{Say it:} \emph{bhedbhāv}
  \item \emph{Say it:} \emph{bhedbhāv}, once more
  \item \emph{Say it:} \emph{āmce śikṣak koṇtāhī bhedbhāv karat nāhīt.} — Our teachers show no discrimination at all.
  \item \emph{Recall:} say the Marathi for a claim, then the Marathi for equality, then say \emph{bhedbhāv} again
\end{itemize}

\begin{tcolorbox}[breakable,colback=teal!4,colframe=teal!35!black,title={Before you move on}]
What does \mr{भेदभाव} mean? (\textquotedblleft{}Discrimination\textquotedblright{}.) Say it once more.
\end{tcolorbox}
\section[\texorpdfstring{\mr{विषमता} (viṣamtā)}{विषमता (viṣamtā)}]{\mr{विषमता} (viṣamtā) — inequality}

\label{lesson:MR-C360-visamta}



\begin{quote}
Before the new one: say the Marathi for equality, then the Marathi for discrimination.
\end{quote}

\subsection*{You'll want to know: \mr{विषमता}}
\textbf{\mr{विषमता}} — \emph{viṣamtā} — \textquotedblleft{}inequality\textquotedblright{}.

\begin{quote}
\mr{श्रीमंत} \mr{आणि} \mr{गरीब} \mr{यांच्यातली} \mr{विषमता} \mr{वाढत} \mr{आहे}.
\end{quote}

\begin{quote}
\emph{śrīmant āṇi garīb yāñcyātlī viṣamtā vāḍhat āhe.}
\end{quote}

\begin{quote}
The inequality between the rich and the poor is growing.
\end{quote}

Feminine.

\mr{विषम}, unequal, + -\mr{ता}; the opposite is \mr{समानता}.

\begin{cousinweb}
\textbf{\mr{विषमता}} (\emph{viṣamtā}) is a learned noun built from Sanskrit \emph{viṣama}, \textquotedblleft{}unequal, uneven\textquotedblright{} (CDIAL 11972), with \emph{-tā}, \textquotedblleft{}-ness\textquotedblright{}.
\end{cousinweb}

\begin{grammarlens}[title={What you've built}]
One more word for giving reasons and opinions.
\end{grammarlens}

\subsection*{Your turn}
\begin{itemize}
\raggedright
  \item \emph{Say it:} \emph{viṣamtā}
  \item \emph{Say it:} \emph{viṣamtā}, once more
  \item \emph{Say it:} \emph{śrīmant āṇi garīb yāñcyātlī viṣamtā vāḍhat āhe.} — The inequality between the rich and the poor is growing.
  \item \emph{Recall:} say the Marathi for equality, then the Marathi for discrimination, then say \emph{viṣamtā} again
\end{itemize}

\begin{tcolorbox}[breakable,colback=teal!4,colframe=teal!35!black,title={Before you move on}]
What does \mr{विषमता} mean? (\textquotedblleft{}Inequality\textquotedblright{}.) Say it once more.
\end{tcolorbox}
\section[\texorpdfstring{\mr{दारिद्र्य} (dāridrya)}{दारिद्र्य (dāridrya)}]{\mr{दारिद्र्य} (dāridrya) — poverty}

\label{lesson:MR-C360-daridrya}



\begin{quote}
Before the new one: say the Marathi for discrimination, then the Marathi for inequality.
\end{quote}

\subsection*{You'll want to know: \mr{दारिद्र्य}}
\textbf{\mr{दारिद्र्य}} — \emph{dāridrya} — \textquotedblleft{}poverty\textquotedblright{}.

\begin{quote}
\mr{शिक्षणामुळे} \mr{अनेक} \mr{कुटुंबं} \mr{दारिद्र्यातून} \mr{बाहेर} \mr{आली}.
\end{quote}

\begin{quote}
\emph{śikṣaṇāmuḷe anek kuṭumbaṁ dāridryātūn bāher ālī.}
\end{quote}

\begin{quote}
Through education many families came out of poverty.
\end{quote}

Neuter.

Neuter and formal; in speech \mr{गरिबी}, from \mr{गरीब}, poor.

\begin{cousinweb}
\textbf{\mr{दारिद्र्य}} (\emph{dāridrya}) is a tatsama, taken from Sanskrit \emph{dāridrya}, \textquotedblleft{}poverty\textquotedblright{} (CDIAL 6297), built on \emph{daridra}, \textquotedblleft{}poor\textquotedblright{}.
\end{cousinweb}

\begin{grammarlens}[title={What you've built}]
One more word for giving reasons and opinions.
\end{grammarlens}

\subsection*{Your turn}
\begin{itemize}
\raggedright
  \item \emph{Say it:} \emph{dāridrya}
  \item \emph{Say it:} \emph{dāridrya}, once more
  \item \emph{Say it:} \emph{śikṣaṇāmuḷe anek kuṭumbaṁ dāridryātūn bāher ālī.} — Through education many families came out of poverty.
  \item \emph{Recall:} say the Marathi for discrimination, then the Marathi for inequality, then say \emph{dāridrya} again
\end{itemize}

\begin{tcolorbox}[breakable,colback=teal!4,colframe=teal!35!black,title={Before you move on}]
What does \mr{दारिद्र्य} mean? (\textquotedblleft{}Poverty\textquotedblright{}.) Say it once more.
\end{tcolorbox}
\section[\texorpdfstring{\mr{साक्षरता} (sākṣartā)}{साक्षरता (sākṣartā)}]{\mr{साक्षरता} (sākṣartā) — literacy}

\label{lesson:MR-C360-saksarta}



\begin{quote}
Before the new one: say the Marathi for inequality, then the Marathi for poverty.
\end{quote}

\subsection*{You'll want to know: \mr{साक्षरता}}
\textbf{\mr{साक्षरता}} — \emph{sākṣartā} — \textquotedblleft{}literacy\textquotedblright{}.

\begin{quote}
\mr{गावातील} \mr{साक्षरता} \mr{गेल्या} \mr{दहा} \mr{वर्षांत} \mr{वाढली}.
\end{quote}

\begin{quote}
\emph{gāvātīl sākṣartā gelyā dahā varṣānt vāḍhlī.}
\end{quote}

\begin{quote}
Literacy in the village grew in the last ten years.
\end{quote}

Feminine.

\mr{साक्षर}, literate, + -\mr{ता}; the opposite is \mr{निरक्षरता}.

\begin{cousinweb}
\textbf{\mr{साक्षरता}} (\emph{sākṣartā}) is a learned noun built from Sanskrit \emph{sākṣara}, \textquotedblleft{}containing letters, lettered\textquotedblright{} (CDIAL 13320), with \emph{-tā}, \textquotedblleft{}-ness\textquotedblright{}.
\end{cousinweb}

\begin{grammarlens}[title={What you've built}]
One more word for giving reasons and opinions.
\end{grammarlens}

\subsection*{Your turn}
\begin{itemize}
\raggedright
  \item \emph{Say it:} \emph{sākṣartā}
  \item \emph{Say it:} \emph{sākṣartā}, once more
  \item \emph{Say it:} \emph{gāvātīl sākṣartā gelyā dahā varṣānt vāḍhlī.} — Literacy in the village grew in the last ten years.
  \item \emph{Recall:} say the Marathi for inequality, then the Marathi for poverty, then say \emph{sākṣartā} again
\end{itemize}

\begin{tcolorbox}[breakable,colback=teal!4,colframe=teal!35!black,title={Before you move on}]
What does \mr{साक्षरता} mean? (\textquotedblleft{}Literacy\textquotedblright{}.) Say it once more.
\end{tcolorbox}
